\section{Introduction}
Introduce the scientific problem before describing the technical approach.
Explain what is already established, what remains unresolved and why the
unresolved question is consequential. A useful introduction narrows from the
broader context to a specific hypothesis that the study can test. Use citations
to support particular statements rather than placing a long list of references
at the end of a general paragraph.

The citation here demonstrates a numbered reference to a real publication; it
is not evidence for the placeholder text in this template.\cite{shannon1948}
Replace this example with literature relevant to the study. The final paragraph
of the introduction should state the question addressed, summarize the strategy
and identify the main contribution without repeating the abstract.

\section{Results}
\subsection{State the principal finding}
Organize the results around scientific findings. Each subsection should make
one clear point, describe the observation that supports it and explain the
comparison needed to interpret it. Distinguish measured quantities from model
predictions, and provide uncertainty alongside estimates. Table~\ref{tab:design}
illustrates a compact three-line table without fabricated experimental values.

\begin{table}[htbp]
  \caption{Example reporting structure. Replace these labels with the quantities
    and comparisons used in the study.}
  \label{tab:design}
  \centering\small
  \begin{tabular*}{\linewidth}{@{\extracolsep{\fill}}lll@{}}
    \toprule
    Component & Purpose & Report \\
    \midrule
    Reference & Comparison & Baseline \\
    Intervention & Main test & Effect size \\
    Control & Validation & Uncertainty \\
    \bottomrule
  \end{tabular*}
\end{table}

Keep figure legends self-contained: explain the panels, define the symbols and
state the sample size and statistical procedure where applicable. The legend
should allow readers to understand the display without searching through the
methods. Use the text to explain what the evidence means, rather than repeating
every value that is already visible in a table or figure.

\subsection{Test robustness and scope}
Describe the checks that distinguish the proposed explanation from plausible
alternatives. Report sensitivity to experimental conditions, analysis choices
and sampling variation. Negative or inconclusive results should be described
with the same precision as the principal finding. State which conclusions are
supported directly and which remain hypotheses for future work.

\section{Discussion}
Explain how the findings answer the question introduced at the beginning of the
article. Relate the result to earlier work and make clear whether it confirms,
extends or challenges existing understanding. Separate the demonstrated
contribution from possible applications that have not yet been tested.

Discuss limitations in concrete terms. Identify the populations, conditions or
assumptions under which the conclusion may not hold, and explain how these
boundaries affect interpretation. End with the scientific implication of the
work rather than a list of every experiment performed.

\section{Methods}
\subsection{Study design and analysis}
Describe the data, materials and procedures in enough detail to support
reproduction. Specify inclusion criteria, preprocessing, controls and the
analysis plan. Define notation when it first appears. For example, a generic
empirical objective can be written as
\begin{equation}
  \mathcal{L}(\theta)=\frac{1}{N}\sum_{i=1}^{N}
  \ell\bigl(f_{\theta}(x_i),y_i\bigr),
  \label{eq:objective}
\end{equation}
where $N$ is the number of observations, $f_{\theta}$ is a parameterized model
and $\ell$ is a loss function. \Eqref{eq:objective} demonstrates
mathematical typography and a cross-reference; it does not prescribe a method.

\subsection{Reproducibility}
Report software versions, parameter choices and sources of randomness. Explain
how uncertainty was estimated and how repeated measurements were handled.
Provide enough information to distinguish independent observations from
technical repetitions. Move extensive derivations and secondary checks to
supplementary material when this improves the flow of the article.

\section{Data availability}
State where the data can be accessed and describe any access restrictions.
Replace this instruction with a repository identifier or a precise explanation.

\section{Code availability}
Provide the code location and the version associated with the reported results.

\section{Acknowledgements}
Name the funding sources and contributors who do not meet authorship criteria.

\section{Author contributions}
Describe each author's contribution using initials or names consistently.

\section{Competing interests}
Declare relevant interests or state that the authors declare no competing interests.
